\documentclass[11pt]{article}
\usepackage[margin=.8in]{geometry}
\usepackage{amsmath,amssymb,booktabs,longtable,enumitem,xurl,hyperref}
\hypersetup{colorlinks=true,urlcolor=blue}
\setlength{\parindent}{0pt}
\setlength{\parskip}{.5em}
\title{Observer-Feedback Prediction Repair:\\Implementation, Experiments, and Remaining Failures}
\author{Collision Avoidance Research}
\date{October 3, 2026}
\begin{document}
\maketitle
\input{report/OBSERVER_FEEDBACK_REPAIR_20261003/summary.tex}

\textbf{Outcome: partial repair.} The previous noisy-estimator campaign could
not issue its first control in any of fourteen scenarios. The revised
implementation removes that common observer/terminal incompatibility:
all fourteen saved initial conditions issue a control in a fresh no-target
counterfactual. In the complete encounter campaign, however, only
\recoveredCases{} of fourteen cases demonstrate avoidance and nominal recovery;
\unavailableCases{} stop because the controller cannot return a solved input.
There are \collisionCases{} observed collisions in the executed portions.
A stopped run is not a successful collision-avoidance experiment.

\section{The distinction behind the repair}
Receiving another measurement can bound estimation error about the next
estimate. It does not by itself bound prediction error about today's
open-loop trajectory. With true state $x$, estimate $\hat x$, nominal center
$c$, and future estimate-feedback input
$u=\bar u+K(\hat x-c)$, write
\[
 \epsilon=x-\hat x,\qquad d=x-c,\qquad
 d^+=(A+BK)d-BK\epsilon.
\]
This affine identity is the relevant correction to propagating only $Ad$.
The current hold is special: $c_0=\hat x_0$ implies $d_0=\epsilon_0$,
so its input perturbation cancels and $d_1=A_0d_0$.
The implementation respects this cancellation, then carries zonotope
generators through the feedback map without repeatedly boxing the state.
Input and slew reserves use the same generators. The shifted initialization
also uses the stored gains to account for a changed measured initial state;
it remains a search reference, never an independently executable backup.

The finite-horizon brake gain is constrained by available control authority.
For its nominal Riccati row $k_i$, set
\[
 r_i=\|[k_iG_i,-k_iE]\|_1,\quad
 a_i=\max\{0,\min(u_{\max}-\bar u_i,\bar u_i-u_{\min})\},\quad
 k_i\leftarrow\min\{1,a_i/r_i\}\,k_i.
\]
The zero-radius case leaves the gain unchanged. This solves the scalar
problem of retaining the largest fraction of that gain while its input
uncertainty fits the available interval. The state error set is then
propagated with the resulting gain; it is not shrunk to force feasibility.

The online observer-set/feedback-tube distinction is consistent with
K\"ohler, M\"uller, and Allg\"ower's output-feedback MPC framework
\cite{Kohler}. Our local affine implementation does not inherit their
nonlinear theorem automatically.

\section{Implemented changes}
\begin{enumerate}[leftmargin=*,itemsep=3pt]
\item \textbf{Preserve information at the estimator interface.}
The initial radar/GNSS least-squares fit now retains its fitted velocity
magnitude instead of replacing it with a 15 m/s prior. Intersected
position-history information now enters target parameter prediction.
A joint ego generator preserves the correlation between yaw and inferred
body velocity: a yaw error rotates the inferred velocity oppositely, so
inertial-velocity errors nearly cancel. Independent component boxes lost
this cancellation. The scenario adapter supplies a yaw-acceleration bound
derived from the known axle friction capacities, permitting correlated
historical radar transport without supplying the controller any true state.

\item \textbf{Predict directional target uncertainty.}
The constant-$A$, constant-$\beta$ family is unchanged. Analytic sensitivities
in traveled arc, course, and curvature, with a Taylor remainder, bound
support in each collision normal. Longitudinal speed uncertainty therefore
does not automatically become equally large lateral uncertainty.
No future target measurement is invented or used to contract this set.

\item \textbf{Separate the nominal terminal neighborhood from noise.}
The former terminal cone subtracted a large uncertainty reserve from a
microscopic nominal radius, making its right-hand side negative.
For the affine terminal feedback matrix $F$, the replacement constructs
\[
 Z=\frac{W+FW+\cdots+F^{s-1}W}{1-\alpha},\qquad
 F^sW\subseteq\alpha W,\quad \alpha<1.
\]
Then $FZ+W\subseteq Z$. A scaled member of this family contains the endpoint
error set and a nominal mean ellipsoid. Physical state/input/slew domains
are still constrained. A full-dimensional disturbance-box floor uses the
existing nominal core, avoiding enormous scaling in unexcited directions.
A one-dimensional velocity-error regression previously produced a scale
of $4.04\times10^9$; the repaired case is physically admissible and solves.

\item \textbf{Retain feasible current-frame slack budgets.}
When the shifted nonlinear anchor satisfies the current affine constraints,
its required robustness slack supplies a feasible cap and the solver
directly optimizes the same CLF. Otherwise the primary restoration problem
runs. Updated observer information can change this cap: it is a feasible
upper bound, not the optimal PCBF value or a proved cross-frame decrease.
The numerical lexicographic allowance is scaled relative to the primary
value and aligned with the solver's feasibility tolerance.

\item \textbf{Repair local collision and RTI semantics.}
At rectangle overlap, the ordinary distance dual can be zero and produce
a zero-gradient infeasible row. A nonzero feasible dual along a least
penetrated axis supplies a signed escape row. First-hold model agreement
is evaluated for the input actually issued; a remote nonlinear suffix
failure does not masquerade as first-hold disagreement. The remaining
affine horizon constraints are retained. A fresh APF seed is still used
only at initialization or after a failed shift; no new control mode is added.
\end{enumerate}

\section{Explicit experimental limits}
The future posterior enclosure is frozen at its current value for these
experiments. Recurring measurements alone do not prove this envelope.
The terminal invariant family is affine: nonlinear remainders, uniform
gain compatibility, and information-set nesting remain unverified.
The finite RTI suffix is not a certified nonlinear continuation.

With nonzero observer uncertainty, the optimizer now imposes hard nominal
collision geometry and minimizes soft robustness deficits across the full
finite horizon. This removes a hard uncertain-tail obstruction, but positive
deficits do not certify collision avoidance for every state in the uncertainty
set. The configured 0.20 m is a mean-model sampled margin, not an achieved
physical lower bound. The observed minimum in this campaign is
\minimumClearance{} m. The controller explicitly retains
\texttt{robustNonlinearSafetyCertified=false}.
The ideal recursive-feasibility theorem in
\path{controller/OBSERVER_ROBUST_PCBF_THEORY.tex} keeps its assumptions;
the implementation table now identifies these remaining gaps.

\section{Scenario validation}
Seven prescribed encounters are tested at 8 and 15 m/s. The target follows
constant tangential acceleration and sideslip. Radar publication is limited
to 50 m, road constraints are disabled, maximum lateral deviation is 10 m,
and the controller sample period is 0.05 s. The seed is 20261003.
Bounded uniform sensor amplitudes are 0.04 m GNSS position,
0.05 m/s GNSS velocity, 0.03 m/s$^2$ IMU acceleration,
0.0015 rad/s gyro, and 0.04 m radar position. The actual NRMM runtime uses
80 Hz samples and a maximum 0.0025 s integration step.

The driver replays each issued hold independently with tight \texttt{ode45}
integration. Recovery requires one second inside the existing error
tolerances $[0.1\text{ m},1^\circ,0.1\text{ m/s},0.05\text{ m/s},0.01\text{ rad/s}]$
after the prescribed interaction. The 2000-hold observation cap is not a
success deadline. Analytic circular scenarios can meet again; their
baseline interaction window includes later encounters within that cap.
Certificate availability is not a failure condition. Actual collision or
failure to return a solved input is a failure.

\begin{longtable}{lr r r l}
\toprule
Scenario & Speed (m/s) & Executed (s) & Minimum (m) & Outcome \\\midrule
\endhead
\input{report/OBSERVER_FEEDBACK_REPAIR_20261003/scenario-table.tex}
\end{longtable}

Maximum successful-call wall time is \maximumCallMs{} ms, including startup;
the maximum after each scenario's first successful call is
\maximumWarmCallMs{} ms. The longest failed call is \maximumFailedCallMs{} ms.
These are runtime measurements under the machine's actual concurrent load,
not isolated solver benchmarks or a proven 100 ms deadline. Startup
estimator construction is recorded separately and excluded from controller
call timing.

\section{Rejected changes and unresolved design issues}
A fixed small first-input correction cap produced a stable off-path
trajectory at approximately 7.69 m lateral error. On the saved frame,
the optimizer exhausted its steering allowance and returned positive CLF
slack; it could only cancel a shifted countersteer. Removing this cap
restored a decreasing CLF direction. It was not retained.

An enlarged APF envelope, opposite-side seed screening, and a persistent
passing normal did not improve the campaign consistently. The persistent
normal could demand the wrong side after the crossing geometry changed.
These exploratory changes were removed; the overlap escape dual remains.
Simply expanding the estimator sideslip domain from 0.005 to 0.06 rad
changed its gains and exposed substantial curvature oscillation even for
straight targets: only one of fourteen cases recovered. This variant was
also rejected.

On the saved 8 m/s turning-crossing failure at 0.55 s, the shifted plan's
predicted brake-input error radius reached 1.14435, producing the empty
physical interval $[+0.144349,-0.144349]$. Removing only this input reserve
made the offline convex problem solvable; removing terminal or collision
constraints did not. The adopted gain projection repairs this contradiction.
The same saved frame still fails first-hold agreement: its pose discrepancy
is only 3.32 mm, but normalized state and CLF discrepancies are approximately
0.0180 and 0.01255 against their respective 0.01 thresholds. Primary and
secondary first inputs are nearly identical, so damping between them does
not reduce this error. The ensuing APF seed has approximately 2.97 m
rectangle penetration in its predicted suffix and cannot be repaired by
its local collision constraints. This separates the trigger from the
unsuccessful reinitialization.

Additional one-frame ablations removed the generic state-error threshold
and then separately doubled the CLF threshold. They returned more controls
with positive measured clearance over that one hold, but do not establish
complete avoidance or recovery. These threshold changes were not adopted.

Remaining failures need improvements to the local planning model and
target estimates, not another arbitrary uncertainty-radius cap. Some
single-convexification steps still fail first-hold agreement, and collision
rows around a poor APF seed can require more motion than the current local
model permits. Damping toward an inaccurate primary point need not repair
that disagreement. Curved truth trajectories also exceed the current
0.005 rad estimator design domain; this is analyzed as a practical model
mismatch rather than an automatic experimental failure.
Offline ablations and the final failure snapshots are retained with the
experimental artifacts. Their conclusions are local evidence, not a proof
that these scenarios are physically unavoidable.

\section{Verification and reproducibility}
The companion directory records the final scenario metrics, source hashes,
configuration, regression outcomes, algebra checks, and offline findings.
The final-source campaign is isolated from subsequent edits. Rejected
variants are labeled separately. The unit tests include joint-error
containment and cancellation, directional target supports, signed overlap
rows, terminal mean admission, repeated-construction consistency, and
degenerate posterior directions. Code Analyzer findings are compared with
the original revision. No nested solver repository, generated native binary,
or unrelated estimator/theory edit is included in the project change.

\begin{thebibliography}{9}
\bibitem{Kohler} J. K\"ohler, M. A. M\"uller, and F. Allg\"ower,
\emph{Robust output feedback model predictive control using online estimation
bounds}, 2021. \url{https://arxiv.org/abs/2105.03427}.
\end{thebibliography}
\end{document}
